\documentclass[11pt,a4paper]{article}
\usepackage{dsmorgan}

\title{The Reflective Property of the Ellipse}
\subtitle{Two coordinate-free proofs}
\author{Drew M. Schaffner}

\begin{document}
\maketitle

\section{Statement}

\begin{theorem}[Reflective property]\label{thm:main}
  Let $E$ be an ellipse with foci $F_1$, $F_2$. At every point $P \in E$, the
  normal line to $E$ bisects the angle $F_1 P F_2$. Consequently, a ray from $F_1$
  that reflects off $E$ at $P$ passes through $F_2$.
\end{theorem}

\section{Setting}

\begin{definition}[Ellipse]\label{def:ellipse}
  Let $F_1, F_2 \in \R^2$ and $a > c \geq 0$, where $2c = \abs{F_1 - F_2}$. The
  \emph{ellipse} with foci $F_1$, $F_2$ and major axis $2a$ is the level set
  \[
    E = \set{X \in \R^2 : g(X) = 2a}, \qquad
    g(X) = \abs{X - F_1} + \abs{X - F_2}.
  \]
\end{definition}

\begin{definition}[Law of reflection]\label{def:reflection}
  A ray travelling in the unit direction $d$ that meets a curve with unit normal $n$
  leaves in the direction $d' = d - 2 (d \cdot n) n$, the mirror image of $d$ in the
  tangent line.
\end{definition}

The circle is the case $c = 0$, and no point needs separate treatment: unlike a
proof using slopes, nothing goes wrong where the tangent or a focal line is
vertical.

\section{Idea of the proof}

The gradient of a distance $\abs{X - F}$ is the unit vector pointing away from $F$.
So $\nabla g$ is a sum of two unit vectors, and a sum of two unit vectors bisects the
angle between them.

\section{First proof: the gradient}

\begin{lemma}[$E$ is a smooth curve]\label{lem:smooth}
  Neither focus lies on $E$. At every $X \in E$, $g$ is differentiable with
  \[
    \nabla g(X) = u_1 + u_2 \neq 0, \qquad
    u_j = \frac{X - F_j}{\abs{X - F_j}}.
  \]
\end{lemma}

\begin{proof}
  By the triangle inequality, $g(F_1) = g(F_2) = 2c < 2a$, so the foci are not on
  $E$. Away from $F_j$, the map $X \mapsto \abs{X - F_j}$ is differentiable with
  gradient $u_j$. If $u_1 + u_2 = 0$, then $X$ lies on the segment $F_1F_2$, so
  $g(X) = 2c < 2a$ and $X \notin E$.
\end{proof}

By \cref{lem:smooth} and the implicit function theorem, $E$ is a smooth curve near
each of its points, with normal direction $\nabla g = u_1 + u_2$.

\begin{proof}[Proof of \cref{thm:main}]
  Let $P \in E$ and set $u_j$ as in \cref{lem:smooth}, so $u_1$ points from $F_1$
  to $P$ and $u_2$ from $F_2$ to $P$. Since $\abs{u_1} = \abs{u_2} = 1$,
  \[
    u_1 \cdot (u_1 + u_2) = 1 + u_1 \cdot u_2 = u_2 \cdot (u_1 + u_2),
  \]
  so $u_1$ and $u_2$ make equal angles with the normal $u_1 + u_2$. Equivalently,
  $-u_1$ and $-u_2$, which point from $P$ to the foci, make equal angles with the
  inward normal. This is the first claim.

  For the second, the incoming ray travels in direction $d = u_1$. Put
  $s = 1 + u_1 \cdot u_2 > 0$, so that $\abs{u_1 + u_2}^2 = 2s$ and the unit normal is
  $n = (u_1 + u_2)/\sqrt{2s}$. Then $d \cdot n = s / \sqrt{2s}$ and
  \[
    d' = u_1 - 2 \cdot \frac{s}{\sqrt{2s}} \cdot \frac{u_1 + u_2}{\sqrt{2s}}
       = u_1 - (u_1 + u_2) = -u_2 ,
  \]
  which is the direction from $P$ towards $F_2$.
\end{proof}

\section{Second proof: Heron's shortest path}

This proof needs no calculus, only the triangle inequality.

\begin{lemma}[Supporting line]\label{lem:support}
  Let $\ell$ be the tangent line to $E$ at $P$. Then $g(X) > 2a$ for every
  $X \in \ell$ with $X \neq P$.
\end{lemma}

\begin{proof}
  The function $g$ is convex, as a sum of two convex functions. Restricted to $\ell$,
  it is a convex function of one variable whose derivative at $P$ is
  $\nabla g(P) \cdot v = 0$ for a direction $v$ of $\ell$. So $P$ minimises $g$ on
  $\ell$, and $g \geq 2a$ there. If $g(X) = 2a$ for some $X \neq P$ on $\ell$, then
  by convexity $g = 2a$ on the whole segment $PX$. Each $\abs{Y - F_j}$ is convex, so
  each must then be affine along $PX$. A distance $\abs{Y - F_j}$ is affine along a
  segment only if the line through it passes through $F_j$, so $\ell$ would contain
  both foci. But $g(F_1) = 2c < 2a$, contradicting $g \geq 2a$ on $\ell$.
\end{proof}

\begin{proof}[Second proof of \cref{thm:main}]
  Let $F_2'$ be the reflection of $F_2$ in $\ell$, so that
  $\abs{X - F_2} = \abs{X - F_2'}$ for all $X \in \ell$. Both foci satisfy
  $g < 2a$, and the set $\set{g < 2a}$ is convex and does not meet $\ell$ by
  \cref{lem:support}, so $F_1$ and $F_2$ lie on the same side of $\ell$. Hence $F_1$
  and $F_2'$ lie on opposite sides, and the segment $F_1F_2'$ meets $\ell$ at a point
  $X_0$.

  For any $X \in \ell$, the triangle inequality gives
  \[
    g(X) = \abs{X - F_1} + \abs{X - F_2'} \geq \abs{F_1 - F_2'} = g(X_0).
  \]
  So $X_0$ minimises $g$ on $\ell$, and by \cref{lem:support} the only minimiser is
  $P$. Therefore $P$ lies on the segment $F_1F_2'$. The line $F_1P$ continues
  straight to $F_2'$, and reflecting that continuation in $\ell$ sends it to $PF_2$.
  This is exactly the law of reflection at $P$.
\end{proof}

\section{Remarks}

\begin{remark}
  The same computation proves the hyperbola's property. With
  $h(X) = \abs{X - F_1} - \abs{X - F_2}$ we have $\nabla h = u_1 - u_2$, and
  \[
    u_1 \cdot (u_1 - u_2) = 1 - u_1 \cdot u_2 = -u_2 \cdot (u_1 - u_2).
  \]
  So the \emph{tangent} bisects the angle between the focal lines. A ray aimed at
  one focus is reflected towards the other.
\end{remark}

\begin{remark}
  The Analytic Geometry notes prove the theorem with slopes and the tangent-of-angle
  formula. That needs a limiting argument at the four points where a focal line is
  vertical. The gradient proof treats every point the same way.
\end{remark}

\end{document}
